(1) 两个正交本征偏振分量的传播相位为 $\phi_f=2\pi n_fd/\lambda_0$、$\phi_s=2\pi n_sd/\lambda_0$。去掉共同快轴相位后，相对相移 $\delta=2\pi(n_s-n_f)d/\lambda_0$，Jones矩阵为 $J=\mathrm{diag}(1,e^{i\delta})$，它保持电场振幅平方不变。
(2) 将45度线偏振输入作用在四分之一波片矩阵上，得 $\mathbf E_{\rm out}=(1,i)/\sqrt2$（当 $\delta=\pi/2$）。两垂直分量等幅且相差四分之一周期，随时间电场矢量端点作圆周轨迹，因此成为圆偏振；旋向取决于相位和时间因子约定。
(3) 若半波片相移 $\delta=\pi$，在自身轴坐标中Jones矩阵为 $\mathrm{diag}(1,-1)$。将波片旋转 $\alpha$，通过坐标旋转共轭作用于入射方位为 $\theta$ 的线偏振，可得输出方位 $\theta_{\rm out}=2\alpha-\theta$（模 $\pi$），意味着光轴旋转 $\alpha$ 可使偏振方向转动 $2\alpha$。
(4) 令双折射差 $\Delta n=n_s-n_f\ne0$。若需四分之一波片，最薄正厚度为 $d_{\lambda/4}=\lambda_0/(4|\Delta n|)$；半波片则为 $d_{\lambda/2}=\lambda_0/(2|\Delta n|)$。当 $\Delta n\to0$，固定厚度没有相位延迟，要达到相同延迟所需厚度发散。
